% !PRINTED 296
\weberpage{296}
%<<CONT
\gesperrt{2. If an irreducible equation, in whose degree more than one prime
divides, splits into factors by successive adjunction of radicals, then a
splitting into $s$ factors of degree $r$ is brought about by adjunction of the
roots of a metacyclic equation of degree $r$.}

This theorem contains as a special case the theorem due to \gesperrt{Abel},
that an irreducible equation whose degree $m$ is not the power of a prime can
be solvable by radicals only if it splits into factors by the adjunction of
the roots of a solvable equation of lower degree, whose degree is a divisor of
$m$\textsuperscript{1)}. With this the question of the solution of an equation
by radicals, or even of the mere reduction of an equation by radicals, is
extraordinarily simplified. One can in fact restrict the further investigation
to equations whose degree is a prime or a prime power, because all other cases
are reduced to these by repeated application of theorem 2.

Thus, for example, the question of all irreducible equations of degree $6$
solvable by radicals admits a downright trivial answer:

\gesperrt{In order to obtain all metacyclic equations of degree $6$ in any field
$\Omega$, adjoin to the field $\Omega$ a square root and form in the extended
field all cubic equations, or adjoin the root of a cubic equation and form in
the extended field all quadratic equations.}

As an example of such an equation we cite the one treated by \gesperrt{Hesse},
on which depend the circular sections of a surface of degree $2$ not
referred to its principal axes. This problem is solved by square roots, if the
principal axes have first been determined by the solution of a cubic
equation\textsuperscript{2)}.

\noindent\rule{4cm}{0.4pt}\par
\textsuperscript{1)} \gesperrt{Abel} gives the theorem without proof in the
memoir ``Sur la rés. algébr. des équations''. Oeuvres complètes
1881, Bd. II, S. 217. Cf. also the memoir by \gesperrt{Galois} in
\gesperrt{Liouville's} Journal, Bd. 11.\par
\textsuperscript{2)} \gesperrt{Hesse}, ``Ueber die Auflösung derjenigen
Gleichungen $6^{\mathrm{ten}}$ Grades etc.'' Crelle's Journal, Bd. 41
(1851).
